\section{RÉSZLETES RENDSZERTERV}\label{ruxe9szletes-rendszerterv}

\subsection{Rendszer architektúra}\label{rendszer-architektuxfara}

A szoftver belső működése egy szigorúan szekvenciális pipeline-ra épül. Az alábbiakban a rendszer szoftveres
komponenseinek technikai kapcsolatait és az adatok transzformációját
ismertetem.

\subsubsection{Adatfolyam}\label{adatfolyam}

\begin{enumerate}
\def\labelenumi{\arabic{enumi}.}
\item
  Nyers adatbeolvasás: A folyamat a kamera által rögzített analóg
  optikai jel digitalizálásával kezdődik. A bemenet egy W×H felbontású,
  3 csatornás (RGB/BGR) mátrix.
\item
  Geometriai korrekció: A nyers képkockán a rendszer lencsetorzítás
  mentesítést végez az előzetes kalibrációs paraméterek alapján. Az adat
  típusa itt még változatlanul kép.
\item
  Következtetés: A korrigált kép a detektálás modulhoz kerül. Ezen a
  ponton az adat jellege megváltozik: a kimenet már nem kép, hanem egy
  strukturált objektumlista, amely tartalmazza a detektált kátyúk
  osztályát, konfidencia értékét és a határoló dobozának koordinátáit.
\item
  Transzformáció: A 2D pixelkoordináták a geometriai modulba kerülnek,
  amely előállítja a valós térbeli (X, Y) koordinátákat (méterben).
\item
  Megjelenítés: A metaadatok (dobozok és távolságok) visszaírásra
  kerülnek az eredeti videofolyamra, létrehozva a felhasználó számára
  látható kimenetet.
\end{enumerate}

\subsection{Fejlesztői környezet}\label{fejlesztux151i-kuxf6rnyezet}

\subsubsection{Fejlesztői munkaállomás}\label{fejlesztux151i-munkauxe1llomuxe1s}

A szoftver fejlesztése, a tanító adathalmazok kezelése és a mélytanulási
modell betanítása, finomhangolása egy dedikált asztali konfiguráción történik. A gép
specifikációjának kiválasztásakor a legfontosabb szempont a GPU
memóriamérete (VRAM) és a CUDA magok száma volt a párhuzamosított
számítások gyorsítása érdekében.

hardver specifikációja: 

\begin{itemize}
\item
  CPU: Intel Core i5-6600 (4 mag, 3.30 GHz), felelős az operációs
  rendszer futtatásáért és az adatok előkészítéséért (data loading) a
  GPU számára.
\item
  RAM: 16GB DDR4, ez a memóriaigény biztosítja a nagyméretű
  képállományok és a tanítási batch-ek akadálymentes betöltését.
\item
  GPU: NVIDIA GeForce GTX 1060 6GB, a konfiguráció kulcseleme. A 6GB
  dedikált videomemória lehetővé teszi a közepes komplexitású
  konvolúciós hálózatok hatékony tanítását és a valós idejű inferencia
  tesztelését magas képfrissítési ráta mellett.
\end{itemize}

\subsubsection{Szimulációs környezet}\label{szoftver-technoluxf3giuxe1k}

A szakirodalomban elterjedt platformok, különösen a CARLA és a BeamNG.tech, összehasonlítása alapján a projekt megvalósításához az utóbbi környezet alkalmazása indokolt számomra. A döntést a vizuális és dinamikai valósághűség, a validálhatóság.

A monokuláris kamerakép alapján történő 2D-3D geometriai képtranszformáció egyik legkritikusabb paramétere a kamera pozíciója és orientációja az útfelülethez képest. Valós körülmények között az út egyenetlenségei, valamint a jármű mozgásállapot változásai folyamatosan módosítják a karosszéria dőlésszögét. Míg a CARLA alapértelmezett  motorja hajlamos túlságosan stabil, a fizikai zajoktól mentes kameraképpel operálni, addig a BeamNG.tech saját fejlesztésű, soft-body fizikája kiemelkedő pontossággal szimulálja a jármű viselkedését a feladatom szempontjából.

\subsubsection{Fejlesztői környezet, programozási nyelv, könyvtárak és eszközök}

A rendszerfejlesztés elsődleges programozási nyelve a Python. Ez a nyelv a jelenlegi iparági szabvány a mesterséges intelligencia és a gépi látás területén. Mégis a fő indok a Python mellett az, hogy a BeamNG.tech szimulátor hivatalos kommunikációs interfésze (BeamNGpy) is Python alapú, ami, így garantálja a komponensek zökkenőmentes integrációját a szimulátorba. Továbbá a projekt kulcstechnológiái (PyTorch, OpenCV, Ultralytics YOLO) natív Python támogatással rendelkeznek. 

Alkalmazott könyvtárak és keretrendszerek:

BeamNGpy: A rendszer egyik legfontosabb specifikus könyvtára. Ez a nyílt forráskódú Python API biztosítja a kétirányú kommunikációt a BeamNG.tech szimulátorral. Ezen keresztül történik a virtuális környezet inicializálása, a vizsgált jármű (ego-vehicle) monokuláris kamerájának paraméterezése, a képi adatfolyamok kinyerése, valamint a validáláshoz szükséges fizikai referenciaadatok (ground truth metrikus pozíciók, járműorientáció) valós idejű lekérdezése.

OpenCV (Open Source Computer Vision Library): A képfeldolgozási pipeline váza. Nem csupán egy hagyományos videofolyam, hanem a BeamNGpy által másodpercenként szolgáltatott egyedi képkockák beolvasásáért és előfeldolgozásáért felel. Továbbá ez a könyvtár biztosítja a laterális távolságméréshez szükséges geometriai és perspektív transzformációk végrehajtását, valamint a detektált kátyúk és a számított metrikus adatok képernyőn történő valós idejű vizuális megjelenítését.

NumPy (Numerical Python): A rendszer matematikai motorja. A memóriafolytonos tárolást alkalmazó NumPy tömbök (ndarray) nagyságrendekkel gyorsabb adatkezelést tesznek lehetővé. A projektben a kamerából származó pixeladatok mátrixként való kezelésén túl a BeamNG.tech-ből érkező téri adatok feldolgozására, valamint a 2D-3D lokalizációhoz elengedhetetlen lineáris algebrai műveletek (pl. homográfia mátrixszal való szorzás, vektorizált koordináta-transzformációk) gyors elvégzésére szolgál.

Ultralytics YOLO (v10) és PyTorch: A kátyúdetektálást végző mesterséges intelligencia modul magja. A PyTorch egy alacsony szintű mélytanulási keretrendszer, amely dinamikus számítási gráfokat használ, míg az Ultralytics könyvtár egy magas szintű absztrakciót biztosít a YOLO (You Only Look Once) architektúra tanításához, finomhangolásához és futtatásához. Ez a modul felel azért, hogy a szimulált képeken a kátyúk 2D-s képkoordinátáit a képtranszformációs fázis számára azonosítsa.

Fejlesztői eszközök:

\begin{itemize}
	\item
	IDE: Visual Studio Code (könnyűsúlyú kódolás és a Python környezetek egyszerű integrációja).
	\item
	Verziókezelés: Git (a forráskód változásainak követésére).
	\item
	Csomagkezelés: Pip / Anaconda (a függőségek és virtuális környezetek
	izolált kezelésére).
\end{itemize}


\subsection{Modulok felépítése}

\subsubsection{Adatmodell és struktúrák}

\paragraph{Kátyú objektum főbb attribútumai}

A rendszer központi adategysége a Kátyú objektum. Amikor a detektor
felismer egy úthibát, példányosít egy ilyen objektumot, amely a
feldolgozási lánc végéig kíséri az adott hibát, folyamatosan bővülve új
attribútumokkal.

\begin{itemize}
\item
  Azonosító (ID): Egyedi egész szám, amely lehetővé teszi a kátyú
  követését több képkockán keresztül.
\item
  Határoló doboz (bounding box): A detektált kátyú bounding box
  koordinátái a képen (x, y, w, h) formátumban.
\item
  Megbízhatóság (confidence): A mélytanuláson alapuló algoritmus által
  visszaadott valószínűségi 0,0 -- 1,0 közötti érték, amely jelzi,
  mennyire biztos a rendszer a detektálás helyességében.
\item
  Valóspozíció (real world coordinates): A lokalizációs modul által
  számított 3D pozícióvektor X\_lat, Y\_long, amely méterben vagy
  centiméterben adja meg a hiba távolságát a jármű referenciapontjához
  képest.
\end{itemize}

\paragraph{Kamera objektum főbb attribútumai}

Ez az osztály felelős az autóra rögzített kamera elhelyezéséének pontos
paramétereiért.

\begin{itemize}
\item
  Magassága (H\_cam): 1,45 méter a földtől.
\item
  Oldalirányú eltolása (T\textsubscript{\_}x): a jármű középvonalától
  0,0 méterre.
\item
  Hosszirányú eltolása (T\_y): a lökhárítótól való távolság 0,85
  méterre.
\item
  A kamera bólintási szöge (0\_pitch): 0,05236 radiánban (3fok).
\end{itemize}

\subsubsection{Fő vezérlési ciklus}

A rendszer működését egy központi vezérlőmodul koordinálja. Ennek a
modulnak a feladata a rendszerindítás, az erőforrások menedzselése, a
feldolgozási lánc (pipeline) elemeinek szinkronizált meghívása, valamint
a futás felügyelete. Programozástechnikai szempontból ez a modul
tartalmazza a belépési pontot (main függvény) és a végtelen ciklust
(while loop).

\paragraph{Inicializálás és erőforrás foglalás}

A program indításakor a vezérlő első feladata a szükséges statikus
erőforrások betöltése a memóriába, még a feldolgozási ciklus megkezdése
előtt. Ekkor történik a konfiguráció betöltése. Mint a rendszer
paramétereinek (pl. kamerakalibrációs mátrixok, detektálási
küszöbértékek, ROI beállítások) beolvasása külső fájlokból. Modell
inicializálás: A YOLO hálózat súlyainak betöltése a GPU memóriájába,
hogy az első képkocka feldolgozása már ne szenvedjen késedelmet.
Kamerakapcsolat létesítés a cv2.VideoCapture interfész megnyitása és a
kamera paramétereinek (felbontás, FPS) ellenőrzése.

\paragraph{A valós idejű feldolgozási ciklus}

A rendszer egy végtelen ciklus működteti, amely addig fut, amíg a
felhasználó leállítási parancsot nem ad. A cikluson belül szigorú
szekvenciális végrehajtás történik, ezzel biztosítva az adatok
konzisztenciáját.

A vezérlő lekéri a legfrissebb nyers képkockát a kamera pufferéből.
Amennyiben az olvasás sikertelen, a modul hibaüzenetet generál vagy
megkísérli az újra csatlakozást.

Sikeres lekérés után a modulhívások láncolata következik. A vezérlő
sorrendben átadja az adatokat az alrendszereknek:

\begin{enumerate}
\def\labelenumi{\arabic{enumi}.}
\item
  Előfeldolgozó: a nyers képből torzításmentesített képet készít.
\item
  Detektor: A bementi képből elkészíti az objektum listát.
\item
  Lokalizáló: Elvégzi konverziót és kinyeri a távolság adatokat.
\item
  Megjelenítő: A feldolgozás során kinyert adatokat rá illeszti a videó
  képre és megjeleníti azt.
\end{enumerate}

Teljesítmény monitorozás: A vezérlő minden ciklusban méri a futási időt,
és ebből valós időben számítja az FPS értéket. Ez az adat
visszacsatolást ad a rendszer terheltségéről.

A szoftver futásának ellenőrzött, manuális megszakítása a felhasználó
által kezdeményezhető az erre a célra dedikált ESC billentyű
lenyomásával. A vezérlő modul a fő feldolgozási ciklus minden
iterációjában figyeli a billentyűzet-eseményeket, és a megszakítási
kérelem észlelésekor azonnal kilép a végrehajtási hurokból, elindítva a
szabályos leállítási mechanizmust.

\subsubsection{Előfeldolgozás}

Képarány helyes átméretezését a korrigált képet a YOLO által megkövetelt fix felbontásra
640x640 pixel méretezi át (cv2.resize). Az eredeti képarány megtartása
mellett skálázza a képet, a fennmaradó üres területeket
(cv2.copyMakeBorder) pedig kitöltő színnel (padding) látja el. Ezzel
elkerülhető a detektált objektumok alakjának torzulása.

Majd a normalizálás és formázás transzformációs lépések a PyTorch keretrendszer bemeneti
követelményeinek teljesítése. A színteret BGR-ről konvertálja RGB-re, a
pixelértékeket a {[}0, 1{]} tartományba normalizálja, végül a mátrixot a
GPU-feldolgozáshoz szükséges tenzor formátummá alakítja.

A kátyú detektálását az Ultralytics YOLOv10 architektúrája végzi. A modell kiválasztását a  \ref{rendelkezesre-allo-nyilt-adathalmazok-es-versenyek-a-temaban}. pontban (\pageref{rendelkezesre-allo-nyilt-adathalmazok-es-versenyek-a-temaban}. oldalon) tárgyalt Optimized Road Damage Detection Challenge, nemzetközi verseny hivatalos rangsora és eredményei indokolják. A megmérettetésen az első helyezést elérő FRDC-SH csapat megoldása a YOLOv10 modellt alkalmazta, amellyel kiemelkedő egyensúlyt teremtettek a sebesség és a detektálási pontosság között.

A hálózat betanítása során alkalmazott módszertan a transfer learning-ra épül. Jelentős fejlesztési előnyt biztosít, hogy a győztes FRDC-SH csapat a GitHub platformon publikálta a előtanított modelljük súlyfájljait.  

Ennek köszönhetően a tanítási folyamatot nem az alapoktól kell elkezdeni. A rendszer betölti a nagy kiterjedésű, nemzetközi adathalmazokon előzetesen betanított és optimalizált súlyokat. Ezt a folyamatot kizárólag egy célzott finomhangolási fázis követi, amelyhez a dedikált BeamNG.tech szimulációs környezetből kinyert, 2D-s virtuális kameraképeket használom fel. Ez a megközelítés garantálja, hogy a modell a szimulátor specifikus vizuális és dinamikai sajátosságaira optimalizálódjon, miközben az előtanított súlyok révén drasztikusan csökken a szükséges tanítási idő és a hardveres erőforrásigény.

\subsubsection{Kátyú lokalizálás}

\paragraph{Számítás menete}

A számítás megkezdése elött megvizsgálja a kátyú osztály példánynak a
megbízhatóság (confidence) értékét, ha az nagyobb mint 0,5 akkor
belekezd a számításba.

A Kátyú objektumból (Bemenet):

\begin{itemize}
\item
  x, y, w, h: A határoló doboz paraméterei pixelben.
\end{itemize}

A Kamera objektumból (Konstansok):

\begin{itemize}
\item
  H\_cam: Kamera magassága (méterben, pl. 1.45).
\item
  T\_x: Kamera oldalirányú eltolása a jármű középpontjától (méterben).
\item
  T\textsubscript{\_}z: Kamera hosszirányú eltolása a jármű
  referenciapontjától (pl. lökhárítótól) (méterben).
\item
  0\_pitch: A kamera bólintási szöge
\end{itemize}

Kamera belső paraméterek:

\begin{itemize}
\item
  f\_x, f\_y: Fókusztávolság pixelben.
\item
  c\_x, c\_y: Az optikai középpont koordinátái pixelben.
\end{itemize}

\begin{enumerate}
\def\labelenumi{\arabic{enumi}.}
\item
  Lépés a talppontot meghatározása. A kátyú térbeli helyzetét annak az
  útfelülethez legközelebb eső pontja határozza meg. A bounding box
  alapján a referenciapont (u\_ref, v\_ref):
\end{enumerate}

\[u\_ ref\  = \ x\  + \ w/2\]

\[v\_ ref\  = \ y\  + \ h\]

\begin{enumerate}
\def\labelenumi{\arabic{enumi}.}
\setcounter{enumi}{1}
\item
  Lépés: a beesési szög számítása. A képpont függőleges koordinátájából
  (v\_ref) meghatározza a fénysugár beesési szögét az optikai tengelyhez
  képest
\end{enumerate}

\[alpha\ \  = \ arctan((v\_ ref\  - \ c\_ y)/f\_ y)\]

A teljes szöget (gamma) az útfelület normálisához (függőlegeshez) képest
a kamera fizikai dőlésszöge (0\_pitch) és a pixeltől függő szög összege
adja:

\[gamma\  = \ 0\_ pitch\  + \ alpha\]

\begin{enumerate}
\def\labelenumi{\arabic{enumi}.}
\setcounter{enumi}{2}
\item
  Lépés: a hosszirányú távolság számítása (Kamera koordinátákban). A
  síkút feltételezés (Z\_worl = 0) alapján egy derékszögű háromszög
  alakul ki, ahol az egyik befogó a kamera magassága (H\_cam), a másik
  pedig a keresett távolság (d\_optical) a talajon.
\end{enumerate}

\[d\_ optical\  = \ H\_ cam\ /\ tan(gamma)\]

Ez az érték a kamera vetületétől mért távolság az úton.

\begin{enumerate}
\def\labelenumi{\arabic{enumi}.}
\setcounter{enumi}{3}
\item
  Lépés: az oldalirányú pozíció számítása (Kamera koordinátákban). A
  vízszintes pixelkoordináta (u\_ref) és a fókusztávolság (f\_x) aránya
  megegyezik a valós oldalirányú kitérés (x\_cam) és a mért távolság
  (d\_optical) arányával:
\end{enumerate}

\[x\_ cam\  = \ (u\_ ref\  - \ c\_ x)\ *\ (d\_ optical\ /\ f\_ x)\]

\begin{enumerate}
\def\labelenumi{\arabic{enumi}.}
\setcounter{enumi}{4}
\item
  Lépés: transzformáció a Jármű-koordinátarendszerbe. Utolsó lépésként a
  kamera koordináta rendszerében kapott értékeket átültetjük a jármű
  rögzített koordinátáiba. Mivel a kamera a szélvédő mögött van a
  hosszirányú távolságból ezt az értéket le kell vonni. Hosszirányú
  pozíció (Y\_long):
\end{enumerate}

\[Y\_ long\  = \ d\_ optical\  - \ T\_ z\]

(Ha Y\_éong negatív, a kátyú már a lökhárító vonala mögött van).

Oldalirányú pozíció (X\_lat):

\[X\_ lat\  = \ x\_ cam\  + \ T\_ x\]

\paragraph{Bólintási zaj
csökkentése}\label{buxf3lintuxe1si-zaj-csuxf6kkentuxe9se}

A számításra legkritikusabb hibaforrása a jármű felfüggesztésének
mozgásából eredő dinamikus bólintó hatás, ami által az állandónak
tekintett pitch (0\_pitch) nem lesz igaz. Így a kátyú számított
hosszirányú pozíciójában jelentős, véletlenszerűnek tűnő ingadozásokat
okoz. Ennek a geometriai zajnak a kiküszöbölésére a rendszer egy
lineáris Kalman szűrőt alkalmazása indokolt, amely a zajos méréseket egy
determinisztikus fizikai mozgásmodellel ötvözi.

A szűrő működésének alapelve azon a fizikai tényen nyugszik, hogy az
úthiba statikus objektum, a jármű pedig tehetetlenségéből adódóan rövid
időintervallumon belül egyenletes sebességgel közeledik felé, így a
kátyú valós távolságának változása nem lehet ugrásszerű.

\subsubsection{Megjelenítés}\label{megjelenuxedtuxe9s}

\paragraph{Grafikus annotáció (Bounding Box
kirajzolás)}\label{grafikus-annotuxe1ciuxf3-bounding-box-kirajzoluxe1s}

A modul bemenetét a detektált és szűrt kátyú objektumok listája képezi.
A rendszer iterál ezen a listán, és minden egyes objektumhoz tartozó
pixelkoordináták (x, y, w, h) alapján egy határoló keretet (bounding
box) rajzol a videóképre. A megvalósításhoz az OpenCV könyvtár rectangle
függvényét alkalmazza a szoftver. A keret színe és vastagsága
konfigurálható, biztosítva a láthatóságot változó fényviszonyok mellett
is.

\paragraph{Metrikus adatok kijelzése}

A puszta detektáláson túl a rendszer legfontosabb kimenete a
távolságbecslés. A modul a boounding box közvetlen közelében, szöveges
formátumban megjeleníti a lokalizációs alrendszer által számított valós
fizikai távolságokat. A szoftver a putText eljárás segítségével írja ki
a jármű orrától mért longitudinális (Y\_long) és a középvonaltól mért
laterális (X\_lat) távolságot méterben vagy centiméterben.

\paragraph{Rendszerállapot és teljesítmény-monitorozás}

A fejlesztés és a valós idejű tesztelés során kritikus fontosságú a
hardveres erőforrások kihasználtságának nyomon követése. A modul a
képernyő sarkában folyamatosan frissülő információs sávot jelenít meg.
Itt kerül kiírásra a vezérlő modul által mért pillanatnyi képfrissítési
sebesség (FPS). Ez az adat szolgál elsődleges validációként arra, hogy a
rendszer képes-e tartani a specifikációban előírt valós idejű működést a
feldolgozás során.